% Section 10: what would extend this work, plus the design for an interactive
% scenario tool. Kept to about one page.
\section{Extending this work}
\label{sec:next}

Better measurement would close the gaps this analysis ran into, the evidence
base now assembled would support analyses this study could not run, and a small
tool would let readers test the findings under their own assumptions.

\subsection*{Additions that would close measurement gaps}

\begin{itemize}
\item \textbf{Add an earnings module to the next Austin music survey.} Five
      questions asked consistently every two years would rebuild the series
      that lapsed after the 2014 Austin Music Census (whose income items
      covered calendar 2013): gross income from music, share of total income
      from music, number of paid performances, whether the respondent holds
      non-music work, and health-coverage source.\footnote{Both Austin music
      censuses are self-selected online instruments (2014 n=3{,}968; 2022
      n=2{,}260), so a restored earnings series would describe respondents
      rather than the Austin music workforce; a probability-sample benchmark
      for that workforce does not exist. Sources: Titan Music Group,
      \emph{The Austin Music Census}, 2015, at
      \url{https://austin.widen.net/content/om7zqn1qbx/pdf/Austin_Music_Census_Interactive_PDF_53115.pdf};
      Sound Music Cities and the City of Austin Music \& Entertainment
      Division, \emph{2022 Greater Austin Music Census}, at
      \url{https://www.austinmusiccensus.org/s/2022-Greater-Austin-Music-Census-SUMMARY-REPORT-v13-optimized.pdf}.}
      Restoring the module would be an adoption rather than a design task,
      because Sound Music Cities already asks income items in other cities
      of the same census cohort (Section~\ref{sec:measurement}).
\item \textbf{Request the administrative records that agencies collect but do
      not publish.} This study stops short in places because the agencies
      holding the records have not released them: applicant demographics for
      Austin's Live Music Fund, the recipient list for the Texas Music
      Incubator Rebate Program at the Governor's office, and the audited
      fund-balance schedule for Austin's cultural funds.\footnote{A requester
      can obtain each of these through a formal public information request
      under the Texas Public Information Act (Texas Government Code
      Chapter 552, at
      \url{https://statutes.capitol.texas.gov/Docs/GV/htm/GV.552.htm}) or the
      City of Austin's equivalent process. This study filed no such requests,
      because the statutory response and Attorney General review windows would
      have delayed publication, and a follow-up round remains the most direct
      way to close these three gaps. Rider 40 to the Governor's office in the
      2026--27 General Appropriations Act requires the office to report
      film-incentive awards recipient by recipient to the Legislative Budget
      Board (the Legislature's fiscal research office) twice a year, and no
      parallel rider covers the music rebate, so the rebate's recipient list
      stays unpublished unless someone asks for it.}
\item \textbf{Track paid performance slots rather than venue names.} Venue
      addresses recycle as one operator replaces another, so counting names
      overstates what closed. Capacity and the number of paid performance
      opportunities per month measure what a working musician actually faces.
\item \textbf{Extend the venue panel to beer-and-wine rooms.} The Texas
      Comptroller's mixed-beverage tax file omits venues without a liquor
      permit, which removes several well-known listening rooms from the panel
      (Cactus Cafe, Meanwhile Brewing, the Carousel Lounge) even when they
      programme live music regularly.\footnote{The Comptroller's mixed-beverage
      gross-receipts file (dataset \texttt{naix-2893} on the Texas Open Data
      Portal, at
      \url{https://data.texas.gov/dataset/Mixed-Beverage-Gross-Receipts/naix-2893})
      covers only holders of a mixed-beverage (liquor) permit. Rooms operating
      on a beer-and-wine permit are reported through a separate Comptroller
      tax file that a follow-up round could join to the same venue list.}
      A join on permit type would recover the beer-and-wine rooms. Rooms with
      no separable alcohol permit, and rooms folded into a hotel permit, sit
      outside any alcohol-tax file and would need a different source, such as
      ticketing or sound-permit records.
\end{itemize}

\subsection*{Analyses this evidence base would support}

\begin{itemize}
\item \textbf{Link venue closures to property records.} Matching this study's
      closure timeline against Travis Central Appraisal District parcel and
      sale records would separate rent pressure from ownership transfer and
      redevelopment, which the news record conflates.\footnote{The Travis
      Central Appraisal District publishes parcel-level records at
      \url{https://www.traviscad.org/}, with owner name, parcel value and
      sale history for every real-property account in Travis County.}
\item \textbf{Follow the cohort rather than the occupation.} Longitudinal
      earnings records, such as the Census Bureau's Longitudinal
      Employer-Household Dynamics program (LEHD, the state unemployment
      insurance records the Bureau assembles into a national worker panel),
      would show whether musicians leave music, leave Austin, or add non-music
      work. Those three outcomes look identical in the cross-sectional
      microdata used here.\footnote{Public LEHD tabulations are at
      \url{https://lehd.ces.census.gov/}; microdata access requires an
      approved Federal Statistical Research Data Center project.}
\item \textbf{Compare Austin against a matched set of metros} on the same
      federal series used here, to separate a music-specific pattern from
      what happens to self-employed workers in any fast-growing metro.
\end{itemize}

\subsection*{A scenario tool for non-technical readers}

This study prints one set of numbers, and readers will want to vary them. A
small interactive tool with the four controls below would let a council
member, a state staffer, or a reporter test how the measured trends play out
under different assumptions.

\begin{center}
\footnotesize
\begin{tabular}{@{}B{2.9cm} P{7.2cm} P{5.5cm}@{}}
\toprule
\textbf{Control} & \textbf{What the user adjusts} & \textbf{What updates on screen} \\
\midrule
Venue revenue trend & The annual real change in core live-music venue receipts,
from continued decline through stabilisation to renewed growth &
Projected count of core live-music rooms and total real receipts to 2032, with
the historical series drawn behind the projection \\
\rowrule
Public support level & The Live Music Fund's share of hotel-tax revenue, and
the split between musician grants and venue grants &
Dollars available, grants per applicant at current award sizes, and the share
of the estimated musician population reached \\
\rowrule
Housing costs & Austin market rent growth, from the recent decline through a
return to the 2015 to 2023 pace &
The market rent index and, separately, hours of work per month of Fair Market
Rent (the U.S. Department of Housing and Urban Development benchmark used to
set voucher payment standards) at the payroll musician wage, which covers
payroll jobs only. The two rent series moved in opposite directions after
2023, as Section~\ref{sec:costs} shows, so the panel would have to display
both \\
\rowrule
State program design & Whether the state rebate stays tied to alcohol tax
remitted or is tied instead to documented artist payments &
Illustrative redistribution across venue sizes, showing which venues gain and
which lose under each rule \\
\bottomrule
\end{tabular}
\end{center}

\vspace{3pt}
{\footnotesize Whoever builds such a tool would keep it honest by drawing
every default from the measured series here, by drawing projections in a
different line style from observed data and labelling them assumptions, and
by saying plainly that the fourth control illustrates a design tradeoff and
predicts no legislative outcome. A static web page with those series embedded
as a small data file is enough for the historical series, which are already
registered in this study's numbers ledger. Two controls would need
code this study does not have: a projection rule behind the venue trend, and
an illustrative redistribution model behind the state-design control, which
has no measured artist-payment data to rest on.
Whether or not anyone builds it, each item in this section would let a later
analyst replace one of this study's inferences with a measurement.}
